% !TeX program = xelatex
\documentclass{resume}
\usepackage{xeCJK}
\usepackage{linespacing_fix}
\usepackage{cite}
\geometry{vmargin=0.43in}

% 随项目提供的 TrueType 字体可在 PDF 中嵌入字形和 Unicode 映射。
\setCJKmainfont[
  Path = fonts/zh_CN-Arphic/,
  AutoFakeBold = 2,
  ItalicFont = gkai00mp.ttf,
  SmallCapsFont = gbsn00lp.ttf,
  SmallCapsFeatures = {FakeBold=2}
]{gbsn00lp.ttf}

\hypersetup{pdftitle={张绍磊 - 投资人版简历}, pdfauthor={张绍磊}, pageanchor=false}

\definecolor{red1}{RGB}{139,0,0}

\definecolor{red2}{RGB}{255,0,0}


\begin{document}
\pagenumbering{gobble}



%***"%"后面的所有内容是注释而非代码，不会输出到最后的PDF中
%***使用本模板，只需要参照输出的PDF，在本文档的相应位置做简单替换即可
%***修改之后，输出更新后的PDF，只需要点击Overleaf中的“Recompile”按钮即可
%**********************************姓名********************************************
\name{张绍磊}

%**********************************联系信息****************************************
% \otherInfo{\faBookmark\ 生源地：北京}{\faPhone\ 13718537922}{\faHome\ \href{https://zhangshaolei1998.github.io/}{个人主页}}{\faEnvelope\ \href{mailto:zhangshaolei20z@ict.ac.cn}{zhangshaolei20z@ict.ac.cn}}
\otherInfo{\faLocationArrow\ 中国人民大学助理教授}{\faBookmark\ 语音交互与 Agentic 双栈研究员}{}{}
\otherInfo{\faPhone\ 13718537922$\;\;$}{\faHome\ \href{https://zhangshaolei1998.github.io/}{个人主页}}{\faEnvelope\ \href{mailto:zhangshaolei98@ruc.edu.cn}{zhangshaolei98@ruc.edu.cn}}{}
% \otherInfo{\faGithub\ \href{https://github.com/zhangshaolei1998}{github.com/zhangshaolei1998}}{}{}{}

%第一个括号里写手机号，第二个写邮箱
%\contactInfo{(+86) 13718537922}{zhangshaolei20z@ict.ac.cn}
%**********************************其他信息****************************************
%在大括号内填写其他信息，最多填写4个，但是如果选择不填信息，
%那么大括号必须空着不写，而不能删除大括号。
%\otherInfo后面的四个大括号里的所有信息都会在一行输出
%如果想要写两行，那就用两次这个指令（\otherInfo{}{}{}{}）即可
%\otherInfo{性别：男}{籍贯：江南}{}{}
%\otherInfo{来历：钱钟书《围城》}{}{}{}
%*********************************照片**********************************************
%照片需要放到images文件夹下，名字必须是you.jpg，如果不需要照片可以不添加此行命令
%0.15的意思是，照片的宽度是页面宽度的0.15倍，调整大小，避免遮挡文字
\yourphoto{0.11}
%**********************************正文**********************************************


%***大标题，下面有横线做分割
%***一般的标题有：教育背景，实习（项目）经历，工作经历，自我评价，求职意向，等等
\section{\textcolor{red1}{\faUser 技术简介}}
\fontsize{10pt}{1ex}\selectfont

\textbf{国内少数同时具备实时语音模型、全模态交互模型和 Agentic 大模型后训练完整研发经验的研究者。}面向下一代实时交互式 Agent，推动大模型从能听、能说、能看，走向能推理、能调用工具、能自主完成任务。

\vspace{2mm}

\textbf{小米大模型核心团队顾问，MiMo-V2-Flash、MiMo-V2.5、MiMo-V2.6 核心贡献者。}现于中国人民大学数据治理实验室开展大模型研究（杜小勇、范举教授领衔）。

\section{\textcolor{red1}{\faBarChart\ 代表成果与产业落地}}
\begin{itemize}[itemsep=2pt,parsep=0.3ex]
    \item \textbf{实时语音模型}。提出 \textbf{\href{https://github.com/ictnlp/StreamSpeech}{StreamSpeech}}，在单个低延时模型中实现语音识别、翻译和合成，相关算法已在华为落地。参与研发实时语音对话模型 \mbox{\textbf{\href{https://github.com/ictnlp/LLaMA-Omni}{LLaMA-Omni}}}，开源项目获 3.1K GitHub 星标。
    \item \textbf{全模态交互模型}。提出 \textbf{\href{https://github.com/ictnlp/Stream-Omni}{Stream-Omni}}，支持文本、视觉、语音的多种组合，实现实时交互。提出 \textbf{\href{https://github.com/ictnlp/LLaVA-Mini}{LLaVA-Mini}}，将单图像处理延时降至 40 ms。
    \item \textbf{Agentic 大模型后训练}。提出 \textbf{\href{https://github.com/ruc-datalab/DeepAnalyze}{DeepAnalyze}}，通过课程式强化学习，使模型自主完成从数据准备到生成分析报告的全流程。项目已在电力、工业等场景落地，获 2 家头部企业感谢信，相关应用累计产生约 \textbf{1000 万元经济收益}。
\end{itemize}

\vspace{1mm}

\textbf{开源与学术影响}。大模型项目累计获得 \textbf{10,000+ GitHub 星标}，模型下载 \textbf{100 万余次}。发表论文 \textbf{40 余篇}，第一或通讯作者论文 \textbf{18 篇}（含 CCF-A 类 11 篇），Google Scholar 引用 \textbf{1500 余次}。


\section{\textcolor{red1}{\faGraduationCap 教育背景}}


%***********一行子标题**************
%***第一个大括号里的内容向左对齐，第二个大括号里的内容向右对齐
%***\textbf{}括号里的字是粗体，\textit{}括号里的字是斜体
\datedsubsection{\textbf{中国科学院计算技术研究所} - \textit{智能信息处理重点实验室\ 博士}}{2020.09 - 2025.06}
\vspace{-0.5mm}
%***********列举*********************
%***可添加多个\item，得到多个列举项，类似的也可以用\textbf{}、\textit{}做强调
\begin{itemize}[parsep=0.3ex]
  \item 导师：冯洋\ 研究员 \quad|\quad 研究方向：大模型、自然语言处理 \quad|\quad 所长特别奖(夏培肃奖)——全所2人 
  % \item \textbf{发表论文}：发表论文\textbf{30余篇}，其中以第一作者发表论文\textbf{15篇}
  % \item \textbf{代表作}：多语言认知模型(百聆1 / 百聆2 / 百聆智能体)、多模态感知模型(StreamSpeech / LLaMA-Omni / LLaVA-Mini)、大模型安全(TruthX / Truth-aware Selection / Auto-RAG)
  % \item \textbf{重大项目}：负责重大工程项目\textbf{4项} \quad|\quad 发明专利\textbf{3项}
  % \item \textbf{国际竞赛}：国际同声传译测评竞赛\textbf{冠军} 
  % \item \textbf{开源项目}：开源项目在GitHub社区累计获得\textbf{5000+星标}，开源模型下载量\textbf{万余次}
\end{itemize}


\datedsubsection{\textbf{北京邮电大学} - \textit{计算机科学与技术实验班\ 本科}}{2016.09 - 2020.06}
\vspace{-0.5mm}
\begin{itemize}[parsep=0.3ex]
  \item 北京市优秀毕业生、国家级奖学金、企业奖学金、北京市三好学生
\end{itemize}


\section{\textcolor{red1}{\faCogs\ 主要研发经历}}


% \textbf{\large \textcolor{red1}{DeepAnalyze：通过Agentic AI自动化实现数据科学全流程闭环}}
% \begin{itemize}[parsep=0.3ex]
%     \item \textbf{DeepAnalyze}：能接入数据库、结构化数据、非结构化数据，自动化完成数据准备、数据分析、可视化、数据建模、数据洞察的全流程数据科学任务，最终生成数据研究报告。
%     % \item 百聆系列研究提供线上服务，支持使用100+种语言进行交互\ [\href{https://arxiv.org/pdf/2306.10968.pdf}{PDF}] [\href{https://github.com/ictnlp/BayLing}{GitHub}] [\href{http://nlp.ict.ac.cn/bayling/demo}{Demo}] [\href{http://nlp.ict.ac.cn/bayling}{Homepage}]
%     \item DeepAnalyze-8B在结构化数据理解、数据科学代码生成、数据分析等任务上优于GPT-4o、Claude等闭源模型。
%     \item 影响：\textcolor{red2}{\textit{首次提出一个通用模型，实现全流程数据科学任务自动化执行。可赋能数据管理、数据分析等垂域场景。}}
% \end{itemize}


% % \qquad\qquad\qquad\quad\textit{\large 数据智能$+$模型智能：具备多语言认知和多模态感知的可信智能系统}

% \vspace{2mm}

% \textbf{\large \textcolor{red1}{多语言认知模型：百聆系列研究，增强LLM语言对齐以实现生成能力的跨语言迁移}}
% \begin{itemize}[parsep=0.3ex]
%     \item \textbf{百聆-1、百聆-2}：通过语言对齐提升LLM多语言能力。支持使用100+语言交互\ [\href{https://arxiv.org/pdf/2306.10968.pdf}{PDF}] [\href{https://github.com/ictnlp/BayLing}{GitHub}] [\href{http://nlp.ict.ac.cn/bayling/demo}{Demo}] [\href{http://nlp.ict.ac.cn/bayling}{Homepage}]
%     \item \textbf{百聆智能体}：增加实时语音交互能力，实现语音翻译、同声传译等功能。
%     % \item 百聆系列研究提供线上服务，支持使用100+种语言进行交互\ [\href{https://arxiv.org/pdf/2306.10968.pdf}{PDF}] [\href{https://github.com/ictnlp/BayLing}{GitHub}] [\href{http://nlp.ict.ac.cn/bayling/demo}{Demo}] [\href{http://nlp.ict.ac.cn/bayling}{Homepage}]
%     \item 在多语言翻译上，百聆-13B达到了GPT-4 95\%的翻译性能，取得开源模型最佳；在通用任务上，百聆-13B达到了GPT-3.5-turbo 89\%的性能；作为基座模型构建翻译、金融、监管等垂直领域LLM。
%     \item 影响：\textcolor{red2}{\textit{首次提出通过语言对齐提升LLM多语言能力的研究；百聆的线上服务被100+高校/企业/单位申请内测。}}
% \end{itemize}



% \textbf{\large \textcolor{red1}{多模态感知模型：从交互方式、语音、视觉多角度构建多模态感知模型}}


% \begin{itemize}[parsep=0.3ex]
%     \item \datedline{\textbf{交互方式}：离线交互$\Rightarrow$实时交互（低延时）}{AAAI 2021/ EMNLP 2021 / EMNLP 2022 / ACL 2022 / ICLR 2023}
%         \begin{itemize}[topsep=-8em,parsep=0.3ex]
%             \item[\footnotesize$\circ$] 实时训练框架：提出“未来知识蒸馏的训练框架” [\href{https://arxiv.org/pdf/2012.12465.pdf}{PDF}]、“实时模型的长度感知框架” [\href{https://aclanthology.org/2022.acl-long.467.pdf}{PDF}]、“鲁棒性实时模型框架” [\href{https://aclanthology.org/2021.autosimtrans-1.1.pdf}{PDF}]；获得国际同声传译竞赛AutoSimTrans 2021冠军。
%             \item[\footnotesize$\circ$] 实时生成策略：提出“基于自监督学习的策略” [\href{https://aclanthology.org/2022.acl-long.176.pdf}{PDF}]、“基于高斯对齐的策略” [\href{https://aclanthology.org/2022.findings-acl.238.pdf}{PDF}]、“基于概率图模型的策略” [\href{https://openreview.net/forum?id=9y0HFvaAYD6}{PDF}] ；滞后2秒，实时性能媲美离线性能。
%             \item[\footnotesize$\circ$] 实时生成能力：提出“MoE多头注意力” [\href{https://aclanthology.org/2021.emnlp-main.581.pdf}{PDF}]、“跨语言信息量化” [\href{https://arxiv.org/pdf/2210.11220.pdf}{PDF}]、“跨模态信息运输” [\href{https://arxiv.org/pdf/2210.12357.pdf}{PDF}] ；显著提升不同语言、模态在实时生成中的表现。
%         \end{itemize}
    
%     \item \datedline{\textbf{语音模型}：“All in One”语音模型同时完成离线/实时语音任务、语音交互任务}{ACL 2024 / Preprint}
%         \begin{itemize}[topsep=-8em,parsep=0.3ex]
%             \item[\footnotesize$\circ$] 统一语音文本理解、生成：可微语音分割DiSeg统一语音文本理解 [\href{https://aclanthology.org/2023.findings-acl.485.pdf}{PDF}]；片段到片段框架Seg2Seg统一语音文本生成 [\href{https://openreview.net/pdf?id=GuErIOGLie}{PDF}] 。
%             \item[\footnotesize$\circ$] 统一离线/实时语音任务：StreamSpeech对齐了语音文本数据两个模态的异构表示，以端到端的方式完成离线和实时（320毫秒延时）的语音识别、语音翻译、语音合成。[\href{https://arxiv.org/pdf/2406.03049.pdf}{PDF}] [\href{https://ictnlp.github.io/StreamSpeech-site/}{Project}]
%             \item[\footnotesize$\circ$] 统一多项语音任务：LLaMA-Omni统一了各种语音任务的数据形式、交互模式，完成多项语音任务 [\href{https://arxiv.org/pdf/2409.06666}{PDF}] 。
%         \end{itemize}
        
%     \item \datedline{\textbf{视觉模型}：高效统一图像视频理解大模型}{Preprint}
%         \begin{itemize}[topsep=-8em,parsep=0.3ex]
%             \item[\footnotesize$\circ$] LLaVA-Mini：高效多模态大模型，将异构的文本数据、图像数据、视频数据统一成token形式来模态预融合。
%             \item[\footnotesize$\circ$] 每张图片推理时间：110毫秒→40毫秒；显存占用：358MB→0.6MB；支持3小时以上的长视频理解。
%         \end{itemize}

%     % \item \datedline{\textbf{端到端实时语音翻译}：滞后3秒，同声传译性能媲美离线翻译}{EMNLP 2022 / ACL 2023 / NeurIPS 2023}
%     %     \begin{itemize}[topsep=-8em,parsep=0.3ex]
%     %         \item[\footnotesize$\circ$] 可微语音分段解决模态间的长度差异：一种无监督语音分段方法，可以将语音分割成以词为单位的片段。可微语音分段可以和任何基座模型一起联合训练，从而得到有益于基座模型的分段结果。
%     %         \item[\footnotesize$\circ$] 跨模态对齐解决模态间的表示差异：将语音到文本的跨模态翻译过程建模为从语音模态到文本模态的信息运输过程。通过动态规划建模所有可能的运输路径，学习语音模态和文本模态间的对齐。
%     %     \end{itemize}
    
%     % \item \datedline{\textbf{级联实时语音翻译}：滞后8个词，实时翻译性能媲美离线翻译}{AAAI 2021 / ACL 2022 / ICLR 2023}
%     %     \begin{itemize}[topsep=-8em,parsep=0.3ex]
%     %         \item[\footnotesize$\circ$] 概率图模型：将实时生成建模为概率图模型，提出了隐马尔科夫Transformer，利用一个隐马模型将生成时刻和生成结果建模为隐藏事件和观测实践，通过优化观测事件的边际似然学习最佳的的生成时刻。
%     %         \item[\footnotesize$\circ$] 基于Agent的实时生成模型：LLM作为生成模块，维护上下文、指令作为memory，引入外部策略模型控制LLM执行生成或者等待action。性能优于Meta发布的SeamlessM4T大模型。[\href{https://aicloud.conestore.cn:30013/bayling/#/}{Demo}]
%     %     \end{itemize}
%     \item 影响：\textcolor{red2}{\textit{⻓期深耕低延时实时多模态模型，系列工作引起来自Meta/FBK/CMU等研究者的关注。代表作StreamSpeech在Twitter社区收获300K+观看，收到国际语音企业Felo/Speechify/To.be合作邀请；开源项目累计收获5K+星标。}}
% \end{itemize}



% \textbf{\large \textcolor{red1}{大模型安全可信：从LLM幻觉、外部攻击、知识匮乏全方位提升模型安全性}}


% \begin{itemize}[parsep=0.3ex]
%     \item \datedline{\textbf{减轻LLM幻觉}：TruthX通过编辑LLM内部表示激活LLM输出真实性，降低幻觉\ [\href{https://arxiv.org/pdf/2402.17811.pdf}{PDF}]}{ACL 2024}
%         \begin{itemize}[topsep=-8em,parsep=0.3ex]
%             % \item[\footnotesize$\circ$] ``LLM在拥有正确知识的情况下，还是有可能产生幻觉'' $\rightarrow$ 如何激活LLM的内部真实性？
%             \item[\footnotesize$\circ$] TruthX在LLM内部表示中构造真实性特征向量，并编辑该特征来激活LLM真实性，从而减轻LLM幻觉。
%             \item[\footnotesize$\circ$] 在TruthfulQA 上，TruthX 将 13 个最先进的LLMs的的真实性平均提高了20\%。TruthX实现了对LLM的精确控制，正向编辑可以激发LLM产生真实回复；而负向编辑完全破坏LLM真实性，生成充满幻觉的回复。
%         \end{itemize}
        
%     \item \datedline{\textbf{防止LLM被恶意攻击}：Truth-Aware Context Selection自适应识别输入中的错误信息\ [\href{https://arxiv.org/abs/2403.07556}{PDF}]}{ACL 2024}
%         \begin{itemize}[topsep=-8em,parsep=0.3ex]
%             \item[\footnotesize$\circ$] 根据真实性对用户输入的上下文进行词级别筛选，丢弃不真实的部分，保护LLM不被错误输入误导。
%             \item[\footnotesize$\circ$] 提升模型对抗恶意攻击、错误输入的能力30\%+
%         \end{itemize}
        
%     \item \datedline{\textbf{数据库检索增强}：Auto-RAG自动化对数据库进行多轮检索，增强LLM知识 [\href{https://arxiv.org/pdf/2411.19443}{PDF}]}{Preprint}
%         \begin{itemize}[topsep=-8em,parsep=0.3ex]
%             \item[\footnotesize$\circ$] Auto-RAG针对复杂场景，实现完全自动化检索，无需人为干预检索轮次和内容。
%             \item[\footnotesize$\circ$] 多跳知识问答等复杂任务上，准确率提升15\%。
%         \end{itemize}
%     \item 影响：\textcolor{red2}{\textit{首个基于LLM可解释性（逐层Attention+FFN）提升LLM安全的研究；被相关Survey收录为代表方法。}}
% \end{itemize}













% \qquad\qquad\qquad\quad\textit{\large 数据智能$+$模型智能：具备多语言认知和多模态感知的可信智能系统}

\vspace{2mm}

\textbf{\large \textcolor{red1}{（1）实时语音与全模态交互}}
\begin{itemize}[parsep=0.3ex]
\item \textbf{实时语音模型}。提出 \textbf{\href{https://arxiv.org/abs/2406.03049}{StreamSpeech}}，将语音识别、翻译与合成统一到同一个低延时语音模型。参与研发 \textbf{\href{https://arxiv.org/abs/2409.06666}{LLaMA-Omni}} 和 \textbf{\href{https://arxiv.org/abs/2505.02625}{LLaMA-Omni 2}}，实现开源大模型的实时语音对话。

\item \textbf{全模态交互与高效视觉理解}。提出 \textbf{\href{https://arxiv.org/pdf/2506.13642}{Stream-Omni}}，支持文本、视觉、语音多种组合下的实时交互。提出 \textbf{\href{https://arxiv.org/abs/2501.03895}{LLaVA-Mini}}，支持图像、高分辨率图像与视频的高效理解，单图像处理延时降至 40 ms，单卡 24 GB GPU 支持超过 1 万帧视频处理。

\item \textbf{时序数据}：围绕“何时读取、何时处理”的实时序列决策问题，提出 \textbf{\href{https://aclanthology.org/2022.acl-long.176/}{Dual Paths}}、\textbf{\href{https://openreview.net/forum?id=9y0HFvaAYD6}{HMT}}、\textbf{\href{https://aclanthology.org/2022.emnlp-main.65/}{ITST}}、\textbf{\href{https://aclanthology.org/2021.findings-emnlp.121/}{GMA}} 等方法，基于对偶学习、隐马尔可夫模型、最优运输与高斯混合模型学习动态对齐策略，在同声传译任务下取得最佳效果。

\item \textbf{\textcolor{red2}{学术影响}}：相关成果形成 18 篇论文，包括 ICLR、NeurIPS、ACL、AAAI 等 CCF-A 类国际会议；LLaVA-Mini被 \textbf{Jiangchuan Liu（IEEE/ACL Fellow、加拿大工程院 Fellow）}评价为高效多模态模型的重要技术参考，被Microsoft的多模态基础模型 Magma 评价为代表性高效视觉语言建模研究；LLaMA-Omni 系列被 \textbf{Shinji Watanabe（IEEE/ISCA Fellow）}与 \textbf{Karen Livescu（ISCA Fellow）}纳入代表性语音大模模型谱系，被\textbf{李海洲（IEEE/ISCA Fellow、新加坡工程院 Fellow）}评价为代表性前沿语音助手；StreamSpeech、HMT 被 Meta AI 评价为语音到语音翻译的重要代表，并被 ByteDance Seed 的技术报告纳入工业级端到端同声传译系统技术体系。

\item \textbf{\textcolor{red2}{开源系统}}：相关开源项目 GitHub 累计获得\textbf{ 6000+ 星标}；自研系统在谷歌、华为、百度举办的 AutoSimTrans \textbf{国际竞赛中获得冠军}，并受邀进行主题报告；相关算法已在华为、小米等公司落地。
\end{itemize}


\textbf{\large \textcolor{red1}{（2）Agentic 大模型后训练与自主执行}}
\begin{itemize}[parsep=0.3ex]
\interlinepenalty=10000
\item \textbf{强化学习与自主执行}。提出 \textbf{\href{https://arxiv.org/abs/2510.16872}{DeepAnalyze}}，通过课程式强化学习训练模型进行多步推理、工具调用和任务执行，完成数据科学全流程与开放式数据研究。提出 \textbf{\href{https://arxiv.org/abs/2602.07371}{DeepPrep}}，以树式推理支持自主数据准备中的多步规划、执行反馈与自我修正。提出 \textbf{\href{https://zhangshaolei1998.github.io/}{DataEvolver}}，利用环境反馈迭代提升大模型训练数据质量。

\item \textbf{记忆方面}：提出 \textbf{\href{https://arxiv.org/abs/2601.06789}{MemGovern}}，长期智能体记忆治理框架，支持记忆写入、筛选、压缩、更新与遗忘，降低记忆污染；提出 \textbf{\href{https://arxiv.org/abs/2601.06966}{RealMem}}，真实交互场景下的智能体记忆评测基准，衡量长期信息的存储、检索、更新与利用能力；提出 \textbf{\href{https://arxiv.org/abs/2411.19443}{Auto-RAG}}，自主迭代检索增强生成，自动完成检索规划、证据选择与答案生成。

\item \textbf{评估方面}：提出 \textbf{\href{https://zhangshaolei1998.github.io/}{CoDA-Bench}}，面向数据密集型代码生成智能体的系统评测基准，覆盖代码生成、数据处理、执行纠错与结果分析，为智能体训练和能力进化提供反馈信号。

\item \textbf{\textcolor{red2}{学术影响}}：相关成果形成 7 篇高水平论文，发表于 ICML、ACL 等 CCF-A 类国际会议；DeepAnalyze 发布后登上 \textbf{HuggingFace 论文日趋榜第一}，模型和数据集下载量最高\textbf{排名 HuggingFace 全站第 6}，国内外社区累计\textbf{浏览量超过 20 万次}；\textbf{Jian Huang（美国统计学会会士、数理统计学会会士）}评价 DeepAnalyze “证明了强化学习训练 agentic LLMs 用于自主数据科学的有效性”。Auto-RAG 被 \textbf{Erik Cambria（IEEE Fellow）}的检索增强基础算法综述收录为代表性方法，并被 \textbf{Aidong Zhang（ACM/IEEE/AIMBE Fellow）}评价为基线方法。

\item \textbf{\textcolor{red2}{开源系统}}：DeepAnalyze 开源项目获得\textbf{ 4.2K GitHub 星标}，线上系统收到国内 300 余所高校与科研机构试用申请；已在电力、工业等行业场景落地，\textbf{获得 2 家头部企业感谢信，累计产生经济收益约 1000 万元}；在部分场景中支撑百万级数据处理与数十万次分析任务，相较传统方式实现 80\%--300\% 效率提升，年节约人力成本数百万元；入选\textbf{第 19 届中国大学生计算机设计大赛大数据主题赛官方支撑智能体}。
\end{itemize}


\textbf{\large \textcolor{red1}{（3）基座模型研发与通用能力优化}}
\begin{itemize}[parsep=0.3ex]
\item \textbf{小米 MiMo 系列}。作为 \textbf{\href{https://arxiv.org/pdf/2601.02780}{MiMo-V2-Flash}}、\textbf{MiMo-V2.5} 和 \textbf{MiMo-V2.6} 的核心贡献者参与模型研发。MiMo-V2-Flash 是 309B 参数 MoE 基座模型，使用 27T tokens 进行预训练。

\item \textbf{后训练阶段}：提出 \textbf{\href{https://arxiv.org/abs/2306.10968}{BayLing-1}} / \textbf{\href{https://arxiv.org/abs/2411.16300}{BayLing-2}}，通过语言对齐实现跨语言迁移，支持 100+ 语言；提出 \textbf{\href{https://aclanthology.org/2025.emnlp-main.328/}{AlignX}}，通过表示对齐提升低资源语言与跨语言迁移能力；提出 \textbf{\href{https://arxiv.org/abs/2508.19839}{PSO-Merging}}，基于粒子群优化低成本融合多专家模型能力。

\item \textbf{推理阶段}：提出 \textbf{\href{https://arxiv.org/abs/2402.17811}{TruthX}}，编辑大模型真实性空间，缓解大模型幻觉；提出 \textbf{\href{https://aclanthology.org/2024.findings-acl.645/}{TACS}}，上下文真实性检测，降低错误输入误导；提出 \textbf{\href{https://aclanthology.org/2025.emnlp-main.537/}{IG-Pruning}}，输入自适应层级剪枝，降低推理成本。

\item \textbf{\textcolor{red2}{学术影响}}：相关成果形成 8 篇论文，包括 ACL、EMNLP 等会议论文及开源技术报告；TruthX 被\textbf{Manuela Veloso（ACM/IEEE Fellow、美国工程院院士）}、\textbf{Anna Korhonen（ACL Fellow）}评价为大模型幻觉缓解代表性工作，被 MIT 研究人员评价其为“\textbf{对大模型内部机制的重要探索}”。推理时的系列研究并被 Google Research、Microsoft Research、NVIDIA 等国际一线研究机构在大模型幻觉机理、大模型推理综述中引用。
\item \textbf{\textcolor{red2}{开源系统}}：BayLing 在中科南京信息高铁研究院线上部署，吸引100余家单位申请内测，被国际开放基准委员会评为“\textbf{Top 100 Open-source Achievements}”；MiMo-V2 系列模型在权威基准\textbf{ AA 榜全球 Top5、国内第一}，并在 OpenRouter 平台取得\textbf{调用量周榜第一}。

\end{itemize}


% \datedsubsection{\textbf{1. \underline{TruthX}}：\textit{通过编辑LLM内部表示来激活LLM输出真实性的推理时对齐方法}}{}


% \begin{itemize}[parsep=0.3ex]
%     \item ``LLM在拥有正确知识的情况下，还是有可能产生幻觉''。TruthX通过在真实空间中编辑LLM的内部表示来激活LLM的真实性，从而减轻LLM的输出幻觉。
%     \item 在TruthfulQA 上，TruthX 将 13 个最先进的LLMs的的真实性平均提高了20\%。TruthX实现了对LLM的精确控制，正向编辑可以激发LLM产生真实回复；而负向编辑完全破坏LLM真实性，生成充满幻觉的回复。
% \end{itemize}

% \datedsubsection{\textbf{2. \underline{Truth-Aware Context Selection}}：\textit{真实性感知的上下文选择方法}}{}
% \begin{itemize}[parsep=0.3ex]
%     \item 根据真实性选择输入context，丢弃上下文中不真实的部分，保护LLM不被带有幻觉的输入/检索结果误导。
% \end{itemize}

% \begin{comment}
    

% \textbf{\large \textcolor{red1}{同声传译：在低延时下（滞后3秒 / 8个词），同声传译性能媲美离线翻译。}}

% \datedsubsection{\textbf{1. \underline{基于受限信息的生成能力}}：\textit{流式模型只能基于当前接收到（受限）的部分源内容生成目标序列}}{}
% \begin{itemize}[parsep=0.3ex]
%     \item \datedline{\textbf{训练方法：未来指导的训练方法}\ [\href{https://arxiv.org/pdf/2012.12465.pdf}{PDF}]}{AAAI 2021 Oral}
%         \begin{itemize}[topsep=-8em,parsep=0.3ex]
%             \item[\footnotesize$\circ$] 提出一种未来指导的训练方式，通过知识蒸馏隐式地将未来信息嵌入到同声传译模型中。有效地使模型获得了隐式预测能力，增强了在信息受限场景下的翻译能力。
%             \item[\footnotesize$\circ$] 首个尝试增强同传模型的预测能力的研究，为后续的研究开辟了新思路。
%         \end{itemize}
        
%     \item \datedline{\textbf{模型框架：基于隐变量的同传框架}\ [\href{https://aclanthology.org/2022.acl-long.467.pdf}{PDF}]}{ACL 2022}
%         \begin{itemize}[topsep=-8em,parsep=0.3ex]
%             \item[\footnotesize$\circ$] 现有的同传框架在信息受限时往往会使同传模型更加重视位置靠前的源端词，即位置偏差。
%             \item[\footnotesize$\circ$] 我们系统分析并论证了位置偏差的影响和产生原因。依此，我们提出了长度感知框架，以将整句长度作为隐变量引入并构造伪整句来减轻同传模型的位置偏差，从而获得更好的同传性能。
%         \end{itemize}
% \end{itemize}

% \datedsubsection{\textbf{2. \underline{同传策略}}：\textit{流式同传模型需要一个策略来权衡翻译质量和延时}}{}
% \begin{itemize}[parsep=0.3ex]
%     % \item \datedline{\textbf{通用策略：基于混合专家的同传策略}\ [\href{https://aclanthology.org/2021.emnlp-main.581.pdf}{PDF}]}{EMNLP 2021 Oral}
%     %     \begin{itemize}[topsep=-8em,parsep=0.3ex]
%     %         \item[\footnotesize$\circ$] 同声传译在不同场景下有不同延时需求。我们将multi-head attention表示为多专家混合（mixture-of-experts）的形式，并提出多专家混合wait-k策略让不同专家学习不同延时下的翻译，以实现通用同传策略。
%     %         \item[\footnotesize$\circ$] 首次提出了通用同传策略的概念，其可以通过一个模型在任意延迟下产生高质量翻译。
%     %     \end{itemize}
        
%     \item \datedline{\textbf{自监督学习：基于自监督学习的同传策略}\ [\href{https://aclanthology.org/2022.acl-long.176.pdf}{PDF}]}{ACL 2022}
%         \begin{itemize}[topsep=-8em,parsep=0.3ex]
%             \item[\footnotesize$\circ$] 同传策略往往缺少明确的约束，影响性能。我们从内部挖掘了两个翻译方向同传策略之间的对偶性约束，并提出一种转置路径的操作以在对偶性约束下联合优化两个翻译方向的同传模型，实现自监督学习。
%             \item[\footnotesize$\circ$] 首个不依赖外部工具或数据，通过在线的自监督学习对同传策略进行监督的研究。
%         \end{itemize}
        
%     % \item \datedline{\textbf{可解释性：基于对齐的同传策略}\ [\href{https://aclanthology.org/2022.findings-acl.238.pdf}{PDF}]}{ACL 2022}
%     %     \begin{itemize}[topsep=-8em,parsep=0.3ex]
%     %         \item[\footnotesize$\circ$] 提出了高斯多头注意力，利用一个高斯先验来联合学习对齐和翻译，然后基于对齐来制定策略。提出的方法无需任何额外的损失函数，就能利用对齐来平衡翻译和延时。
%     %         \item[\footnotesize$\circ$] 首次利用对齐将同声传译中的翻译和延时整合在统一的框架中。
%     %     \end{itemize}

%     \item \datedline{\textbf{概率图模型：基于隐马尔科夫模型的同传策略}\ [\href{https://openreview.net/forum?id=9y0HFvaAYD6}{PDF}]}{ICLR 2023 Spotlight}
%         \begin{itemize}[topsep=-8em,parsep=0.3ex]
%             \item[\footnotesize$\circ$] 提出了隐马尔科夫Transformer，利用一个隐马模型将翻译时刻和翻译结果建模为隐藏事件和观测实践，通过优化观测事件的边际似然学习最佳的的翻译时刻。
%             \item[\footnotesize$\circ$] 首次将同声传译任务建模为概率图模型。取得公开数据集上的SOTA性能，滞后8个词时媲美离线性能。
%         \end{itemize}
% \end{itemize}

% \datedsubsection{\textbf{3. \underline{多模态差异}}：\textit{语音和文本之间的模态鸿沟}}{}
% \begin{itemize}[parsep=0.3ex]
%     \item \datedline{\textbf{模态差异：基于信息运输的同声传译方法}\ [\href{https://arxiv.org/pdf/2210.12357.pdf}{PDF}]}{EMNLP 2022 Oral}
%         \begin{itemize}[topsep=-8em,parsep=0.3ex]
%             \item[\footnotesize$\circ$] 将翻译过程建模为从源端到目标端的信息运输。通过量化运输的信息量，模型在累计信息充分之后开始翻译。
%             \item[\footnotesize$\circ$] 同时适配级联式和端到端语音到文本同传。
%         \end{itemize}
        
%     \item \datedline{\textbf{长度差异：基于可微分段的同声传译方法}}{ACL 2023}
%         \begin{itemize}[topsep=-8em,parsep=0.3ex]
%             \item[\footnotesize$\circ$] 针对语音切分不可导的限制，我们提出了一种可微的分段方法，其可以和任何基座模型一起联合训练，从而得到有益于基座模型的分段结果。
%             \item[\footnotesize$\circ$] 首个可微分的分段方法。取得了公开数据集上的SOTA性能，滞后3秒时媲美离线翻译性能。
%         \end{itemize}
% \end{itemize}
% % \datedsubsection{\textbf{3. \underline{鲁棒的同传模型}}：\textit{同传模型的输出不能被修改，因此面临鲁棒性挑战}}{}
% % \begin{itemize}[parsep=0.3ex]
% %     \item \datedline{\textbf{字符级同传模型}\ [\href{https://aclanthology.org/2021.autosimtrans-1.1.pdf}{PDF}]}{AutoSimTrans@NAACL 2021}
% %         \begin{itemize}[topsep=-8em,parsep=0.3ex]
% %             \item[\footnotesize$\circ$] 真实场景下，文本同传往往是流式语音识别的下游任务，因此输入往往不准确。为此，我们提出一种对于输入更加鲁棒的字符级同传模型。字符级同传模型有效地避免了语音识别的噪声对同传性能的影响
% %             \item[\footnotesize$\circ$] 提出的系统获得国际同声传译测评冠军。在之后测评中，字符级同传模型也成为标准范式。
% %         \end{itemize}
        
% %     \item \datedline{\textbf{信息级同传模型}\ [\href{https://arxiv.org/pdf/2210.11220.pdf}{PDF}]}{EMNLP 2022}
% %         \begin{itemize}[topsep=-8em,parsep=0.3ex]
% %             \item[\footnotesize$\circ$] 现有的同传方法往往在词级别进行决策。这会严重受到源和目标序列长度的影响，在某些语向上性能下降。
% %             \item[\footnotesize$\circ$] 我们提出了信息级同传，量化不同词所包含的信息量，进而在信息层面进行决策。信息级同传能有效地避免源和目标端长度差异和功能词等问题对同传模型的影响，从而进行更鲁棒的决策。
% %         \end{itemize}
% % \end{itemize}
% \end{comment}



% \section{\textcolor{red1}{\faTrophy\ 国际竞赛}}
% \datedsubsection{国际同声传译测评AutoSimTrans 2021流式输入赛道\ \textcolor{red2}{\textbf{冠军}}（由百度、华为、谷歌举办）}{2021.06}
% \datedsubsection{Interdisciplinary Contest in Modeling\  Honorable Mention}{2019.04}


% \newpage

% \section{\textcolor{red1}{\faCogs\ 项目经历}}
% \datedsubsection{\textbf{与华为合作的产学研项目``语音和实时翻译技术的研究项目''} - \textit{主要骨干}}{2019.12 - 2021.12}
% \begin{itemize}[parsep=0.5ex]
%   \item 提出一种未来指导的同传模型训练算法，搭建一套机器同声传译系统，该系统可以完成多个延时下的实时翻译。
%   \item 以第一作者发表论文一篇；作为第一发明人，完成一项发明专利申请；基于自研算法搭建了同传演示系统。
% \end{itemize}

% \datedsubsection{\textbf{计算所创新课题 ``基于神经网络的同声传译方法研究''} - \textit{主要骨干}}{2019.12 - 2021.12}
% \begin{itemize}[parsep=0.5ex]
%   \item 负责研究机器同声传译算法的研究，形成多项学术成果，通过项目验收。
%   \item 基于自研算法的同声传译系统在在由百度、谷歌、华为联合举办的第二届国际机器同声传译比赛流式输入赛道 (AutoSimTrans 2021) 中获得冠军，并受邀于研讨会进行技术分享。
% \end{itemize}

% \section{\textcolor{red1}{\faCogs\ 项目经历}}
% \datedsubsection{\textbf{国家重点研发计划子课题（保密）} - \textit{主要骨干}}{2020.05 - 2022.06}
% \begin{itemize}[parsep=0.3ex]
%   \item 国家级：负责同声传译算法的研究与实现。对提出算法进行系统性测试。最终搭建了算法的演示系统。
%   % \item 完成相关设计报告、测试报告、演示系统等；顺利通过子课题项目验收。
% \end{itemize}

% \datedsubsection{\textbf{国家重点研发计划项目``基于神经网络的汉泰机器翻译研究''} - \textit{主要骨干}}{2021.06 - 2022.07}
% \begin{itemize}[parsep=0.3ex]
%   \item 国家级：负责项目相关算法的整合测试和算法测试报告、结项报告的撰写。
%   % \item 完成相关设计报告、测试报告、演示系统等；顺利通过子课题项目验收。
% \end{itemize}

% \datedsubsection{\textbf{计算所创新课题 ``基于神经网络的同声传译方法研究''} - \textit{主要骨干}}{2019.12 - 2021.12}
% \begin{itemize}[parsep=0.3ex]
%   \item 所级：负责研究机器同声传译算法的研究，形成多项学术成果，取得国际同声传译测评冠军。
% \end{itemize}

% \datedsubsection{\textbf{与华为合作的产学研项目``语音和实时翻译技术的研究项目''} - \textit{主要骨干}}{2019.12 - 2021.12}
% \begin{itemize}[parsep=0.3ex]
%   \item 校企合作：探索同传模型训练算法，搭建一套机器同声传译系统。基于自研算法搭建了同声传译系统。
% \end{itemize}

% \datedsubsection{\textbf{百聆大模型的研究和部署} - \textit{主要骨干}}{2023.04 - 2023.07}
% \begin{itemize}[parsep=0.3ex]
%   \item 所级：负责130亿参数的百聆大模型的构建、训练和线上服务部署。线上服务已获得100+单位申请内测。
% \end{itemize}





\section{\textcolor{red1}{\faCommenting\ 开源项目}}

\begin{itemize}[parsep=0.3ex]

\item \datedline{\textbf{DeepAnalyze} [\href{https://github.com/ruc-datalab/DeepAnalyze}{GitHub}]：自主数据科学大模型}{4.2K星标}

\item \datedline{\textbf{LLaMA-Omni} [\href{https://github.com/ictnlp/LLaMA-Omni}{GitHub}]：实时语音对话大模型}{3.1K星标}

\item \datedline{\textbf{MiMo-V2-Flash} [\href{https://github.com/XiaomiMiMo/MiMo-V2-Flash}{GitHub}]：新一代MoE基座模型}{1.3K星标}

\item \datedline{\textbf{StreamSpeech} [\href{https://github.com/ictnlp/StreamSpeech}{GitHub}]：统一实时语音模型}{1.3K星标}

\item \datedline{\textbf{LLaVA-Mini} [\href{https://github.com/ictnlp/LLaVA-Mini}{GitHub}]：低延时视觉语言模型}{576星标}

\item \datedline{\textbf{Stream-Omni} [\href{https://github.com/ictnlp/Stream-Omni}{GitHub}]：实时三模态交互大模型}{388星标}

\item \datedline{\textbf{BayLing} [\href{https://github.com/ictnlp/BayLing}{GitHub}]：多语言对齐大模型}{314星标}

\item \datedline{\textbf{Auto-RAG} [\href{https://github.com/ictnlp/Auto-RAG}{GitHub}]：自主迭代检索增强}{235星标}

\item \datedline{\textbf{TruthX} [\href{https://github.com/ictnlp/TruthX}{GitHub}]：推理时幻觉缓解}{144星标}

\item \datedline{\textbf{MemGovern} [\href{https://github.com/QuantaAlpha/MemGovern}{GitHub}]：长期记忆治理框架}{126星标}


\item \datedline{\textbf{Awesome Simultaneous Translation} [\href{https://github.com/zhangshaolei1998/Awesome-Simultaneous-Translation}{GitHub}]：同声传译资源库}{580星标}

\item \datedline{\textbf{Awesome Data Agents} [\href{https://github.com/HKUSTDial/awesome-data-agents}{GitHub}]：数据智能体资源库}{568星标}

\end{itemize}







\section{\textcolor{red1}{\faLightbulbO\ 发明专利}}
\datedsubsection{\textbf{发明专利：一种翻译模型的训练方法、翻译方法以及设备} 202011466833.9}{2021}
\datedsubsection{\textbf{发明专利：利用交互式翻译指令的大规模语言模型的指令微调方法} 202310729758.8}{2023}
\datedsubsection{\textbf{发明专利：基于大型多模态模型的视觉表征压缩方法及装置} 202411980108.1}{2024}
\datedsubsection{\textbf{发明专利：一种针对全模态大模型的模态对齐方法及其系统} 202510766248.7}{2025}




\section{\textcolor{red1}{\faCertificate\ 所获荣誉}}

\begin{itemize}[parsep=0.3ex]
%\fontsize{11pt}{0.2ex}\selectfont
    % \item \datedline{北京市三好学生、北京市优秀毕业生}{2018/2020}
    % \item \datedline{中国科学院大学\ 优秀共产党员}{2021}
    % \item \datedline{中国科学院大学\ 三好学生、三好学生标兵}{2021/2022/2023/2024}
    % \item \datedline{中国科学院计算技术研究所\ 一等奖学金}{2022/2023}
    % \item \datedline{国家奖学金（博士）}{2022}
    
    \item \datedline{国际开放基准委员会“Top 100 Open-source Achievements”}{2023}
    \item \datedline{谷歌、华为、百度联合举办的国际竞赛AutoSimTrans2021冠军}{2021}
    \item \datedline{中国科学院计算技术研究所\ 所长特别奖（夏培肃奖）\textcolor{red2}{\textit{计算所最高学生荣誉 全所2人}}}{2022}
    \item \datedline{中国科学院王守武奖学金}{2025}
    \item \datedline{国家奖学金}{2017/2022}
\end{itemize}




\section{\textcolor{red1}{\faStreetView\ 学术服务}}
\begin{itemize}[parsep=0.3ex]
%\fontsize{11pt}{0.2ex}\selectfont
    \item 2020-2024中国中文信息学会青年工作委员会学生执委会\ \textcolor{red2}{\textit{主任}}
    \item CCF-A类国际会议 ACL ARR \textcolor{red2}{\textit{领域主席}} (Area Chair)
    \item 国内会议 YSSNLP 2024、CCL 2024 \textcolor{red2}{\textit{论坛主席}} (Session Chair)
    \item NeurIPS、ACL、EMNLP、ICLR、ICML等国际会议、ACM Computing Surveys期刊 \textcolor{red2}{\textit{审稿人}}
    \item CSSNLP 2020/2021/2023 学生研讨会 \textit{程序委员会主席}
    \item EMNLP 2021/2022、COLING 2022论文预讲会 \textit{程序委员会主席}
    % \item EMNLP 2021、EMNLP 2022、COLING 2022论文预讲会\ \textit{程序委员会主席}
\end{itemize}

\section{\textcolor{red1}{\faBook\ 论文发表}}
\large\textbf{以第一作者/通讯作者在人工智能国际顶级会议发表论文\ \textcolor{red2}{18}\ 篇 （$\dagger$：通讯作者）}
\fontsize{9pt}{9.7pt}\selectfont
\begin{enumerate}[label={[\arabic*]},parsep=0ex]
% \normalsize
    \item \textbf{Shaolei\ Zhang}, Ju Fan, Meihao Fan, Guoliang Li, Xiaoyong Du. DeepAnalyze: Agentic Large Language Models for Autonomous Data Science. \textit{ICML 2026} (\textcolor{red2}{CCF-A}). [\href{https://arxiv.org/pdf/2510.16872}{PDF}] [\href{https://github.com/ruc-datalab/DeepAnalyze}{Code}] [\href{https://ruc-deepanalyze.github.io/}{Homepage}]

    \item Yuxin Zhang, Ju Fan, Meihao Fan, \textbf{Shaolei\ Zhang}$^\dagger$, Xiaoyong Du. CoDA-Bench: Can Code Agents Handle Data-Intensive Tasks? \textit{ICML 2026} (\textcolor{red2}{CCF-A}).

    \item \textbf{Shaolei\ Zhang}, Qingkai Fang, Zhe Yang, Yang Feng. LLaVA-Mini: Efficient Image and Video Large Multimodal Models with One Vision Token. \textit{ICLR 2025} (\textcolor{red2}{CCF-A}). [\href{https://arxiv.org/pdf/2501.03895}{PDF}] [\href{https://github.com/ictnlp/LLaVA-Mini}{Code}]

    \item \textbf{Shaolei\ Zhang}, Tian Yu, Yang Feng. TruthX: Alleviating Hallucinations by Editing Large Language Models in Truthful Space. \textit{ACL 2024} (\textcolor{red2}{CCF-A}). [\href{https://arxiv.org/pdf/2402.17811.pdf}{PDF}] [\href{https://github.com/ictnlp/TruthX}{Code}]

    \item \textbf{Shaolei\ Zhang}, Qingkai Fang, Shoutao Guo, Zhengrui Ma, Min Zhang, Yang Feng. StreamSpeech: Simultaneous Speech-to-Speech Translation with Multi-task Learning. \textit{ACL 2024} (\textcolor{red2}{CCF-A}). [\href{https://arxiv.org/pdf/2406.03049.pdf}{PDF}] [\href{https://github.com/ictnlp/StreamSpeech}{Code}]
     % \item Tian Yu*, \textbf{Shaolei\ Zhang}*, Yang Feng. Truth-Aware Context Selection: Mitigating the Hallucinations of Large Language Models Being Misled by Untruthful Contexts. \textit{Preprint 2024}.[\href{https://arxiv.org/pdf/2403.07556.pdf}{PDF}] [\href{https://github.com/ictnlp/TACS}{Code}]

    \item \textbf{Shaolei\ Zhang}, Yang Feng. Unified Segment-to-Segment Framework for Simultaneous Sequence Generation. \textit{NeurIPS 2023} (\textcolor{red2}{CCF-A}).[\href{https://openreview.net/pdf?id=GuErIOGLie}{PDF}] [\href{https://github.com/ictnlp/Seg2Seg}{Code}]
    
    \item \textbf{Shaolei\ Zhang}, Yang Feng. Hidden Markov Transformer for Simultaneous Machine Translation. \textit{ICLR 2023} (\textcolor{red2}{CCF-A, Spotlight}). [\href{https://openreview.net/forum?id=9y0HFvaAYD6}{PDF}] [\href{https://github.com/ictnlp/HMT}{Code}]

    \item \textbf{Shaolei\ Zhang}, Yang Feng, Liangyou Li. Future-Guided Incremental Transformer for Simultaneous Translation. \textit{AAAI 2021} (\textcolor{red2}{CCF-A, Oral}). [\href{https://arxiv.org/pdf/2012.12465.pdf}{PDF}]
    
    \item \textbf{Shaolei\ Zhang}, Yang Feng. Modeling Dual Read/Write Paths for Simultaneous Machine Translation. \textit{ACL 2022} (\textcolor{red2}{CCF-A}). [\href{https://aclanthology.org/2022.acl-long.176.pdf}{PDF}] [\href{https://github.com/ictnlp/Dual-Path}{Code}]
    
    \item \textbf{Shaolei\ Zhang}, Yang Feng. Reducing Position Bias in Simultaneous Machine Translation with Length-Aware Framework. \textit{ACL 2022} (\textcolor{red2}{CCF-A}). [\href{https://aclanthology.org/2022.acl-long.467.pdf}{PDF}]

    \item \textbf{Shaolei\ Zhang}, Yang Feng. 基于连接时序分类解码器的实时语音翻译方法. \textit{计算机学报 2025} (\textcolor{red2}{CCF-A}).

    \item \textbf{Shaolei\ Zhang}, Yang Feng. 实时翻译研究综述. \textit{中文信息学报 2025} (\textcolor{red2}{CCF-B}).
    

    
    \item \textbf{Shaolei\ Zhang}, Yang Feng. Information-Transport-based Policy for Simultaneous Translation. \textit{EMNLP 2022} (\textcolor{red2}{CCF-B, Oral}). [\href{https://arxiv.org/pdf/2210.12357.pdf}{PDF}] [\href{https://github.com/ictnlp/ITST}{Code}]

    \item \textbf{Shaolei\ Zhang}, Yang Feng. Universal Simultaneous Machine Translation with Mixture-of-Experts Wait-k Policy. \textit{EMNLP 2021} (\textcolor{red2}{CCF-B, Oral}). [\href{https://aclanthology.org/2021.emnlp-main.581.pdf}{PDF}] [\href{https://github.com/ictnlp/MoE-Waitk}{Code}]

    \item \textbf{Shaolei\ Zhang}, Yang Feng. End-to-End Simultaneous Speech Translation with Differentiable Segmentation. \textit{ACL 2023 findings}.[\href{https://aclanthology.org/2023.findings-acl.485.pdf}{PDF}] [\href{https://github.com/ictnlp/DiSeg}{Code}]

    \item \textbf{Shaolei\ Zhang}, Yang Feng. Gaussian Multi-head Attention for Simultaneous Machine Translation. \textit{ACL 2022 findings} . [\href{https://aclanthology.org/2022.findings-acl.238.pdf}{PDF}] [\href{https://github.com/ictnlp/GMA}{Code}]

    \item \textbf{Shaolei\ Zhang}, Shoutao Guo, Yang Feng. Wait-info Policy: Balancing Source and Target at Information Level for Simultaneous Machine Translation. \textit{EMNLP 2022 findings}. [\href{https://arxiv.org/pdf/2210.11220.pdf}{PDF}] [\href{https://github.com/ictnlp/Wait-info}{Code}]
    
    \item \textbf{Shaolei\ Zhang}, Yang Feng. Modeling Concentrated Cross-Attention for Neural Machine Translation with Gaussian Mixture Model. \textit{EMNLP 2021 findings}. [\href{https://aclanthology.org/2021.findings-emnlp.121.pdf}{PDF}]
    
    \item \textbf{Shaolei\ Zhang}, Yang Feng. ICT’s System for AutoSimTrans 2021: Robust Char-Level Simultaneous Translation. \textit{AutoSimTrans@NAACL 2021 workshop}. [\href{https://aclanthology.org/2021.autosimtrans-1.1.pdf}{PDF}]
    

    \item \textbf{Shaolei\ Zhang}, Qingkai Fang, Zhuocheng Zhang, et al. BayLing: Bridging Cross-lingual Alignment and Instruction Following through Interactive Translation for Large Language Models. \textit{Technical Report 2023}. [\href{https://arxiv.org/pdf/2306.10968.pdf}{PDF}] [\href{https://github.com/ictnlp/BayLing}{Code}]

    \item \textbf{Shaolei\ Zhang}, Kehao Zhang, Qingkai Fang, Shoutao Guo, Yan Zhou, Xiaodong Liu, Yang Feng. BayLing 2: A Multilingual Large Language Model with Efficient Language Alignment. \textit{Technical Report 2024}. [\href{https://arxiv.org/pdf/2411.16300}{PDF}] [\href{https://github.com/ictnlp/BayLing}{Code}]

    \item \textbf{Shaolei\ Zhang}, Shoutao Guo, Qingkai Fang, Yan Zhou, Yang Feng. Stream-Omni: Simultaneous Multimodal Interactions with Large Language-Vision-Speech Model. \textit{Preprint 2025}. [\href{https://arxiv.org/pdf/2506.13642?}{PDF}][\href{https://github.com/ictnlp/Stream-Omni}{Code}]

    \item Qihao Wang, Ziming Cheng, Shuo Zhang, Fan Liu, Rui Xu, Heng Lian, Kunyi Wang, Xiaomin Yu, Jianghao Yin, Sen Hu, Yue Hu, \textbf{Shaolei\ Zhang}$^\dagger$, Yanbing Liu$^\dagger$, Ronghao Chen$^\dagger$, Huacan Wang$^\dagger$. MemGovern: Enhancing Code Agents through Learning from Governed Human Experiences. \textit{Preprint 2026}. [\href{https://arxiv.org/pdf/2601.06789}{PDF}] [\href{https://github.com/QuantaAlpha/MemGovern}{Code}]


    
    
    
    % \item \textbf{Shaolei\ Zhang}, Gang Lu, Kai Shuang. Opinion Knowledge Injection Network for Aspect Extraction. \textit{ICONIP 2019} (\textcolor{red2}{CCF-C, Oral}). [\href{https://link.springer.com/chapter/10.1007/978-3-030-36711-4_56}{PDF}]

    % \item Shoutao Guo, \textbf{Shaolei\ Zhang}, Zhengrui Ma, Min Zhang, Yang Feng. SiLLM: Large Language Models for Simultaneous Machine Translation. \textit{Preprint 2024}. [\href{https://arxiv.org/pdf/2402.13036.pdf}{PDF}] [\href{https://github.com/ictnlp/SiLLM}{Code}]

    % \item Shoutao Guo, \textbf{Shaolei\ Zhang}, Yang Feng. Glancing Future for Simultaneous Machine Translation. \textit{ICASSP 2024} (\textcolor{red2}{CCF-B Oral}). [\href{https://arxiv.org/pdf/2309.06179.pdf}{PDF}] [\href{https://github.com/ictnlp/Glance-SiMT}{Code}]

    % \item Zhengrui Ma, \textbf{Shaolei\ Zhang}, Shoutao Guo, Chenze Shao, Min Zhang, Yang Feng. Non-autoregressive Streaming Transformer for Simultaneous Translation. \textit{EMNLP 2022} (\textcolor{red2}{CCF-B Oral}). [\href{https://aclanthology.org/2023.emnlp-main.314.pdf}{PDF}] [\href{https://github.com/ictnlp/NAST}{Code}]

    % \item Shoutao Guo, \textbf{Shaolei\ Zhang}, Yang Feng. Learning Optimal Policy for Simultaneous Machine Translation via Binary Search. \textit{ACL 2023} (\textcolor{red2}{CCF-A}).
    
    % \item Shoutao Guo, \textbf{Shaolei\ Zhang}, Yang Feng. Turning Fixed to Adaptive: Integrating Post-Evaluation into Simultaneous Machine Translation. \textit{EMNLP 2022 findings} (\textcolor{red2}{CCF-B}). [\href{https://arxiv.org/pdf/2210.11900.pdf}{PDF}] [\href{https://github.com/ictnlp/PED-SiMT}{Code}]
    
\end{enumerate}

\large\textbf{合作论文\ \textcolor{red1}{21}\ 篇}
\fontsize{9pt}{9.7pt}\selectfont
\begin{enumerate}[label={[\arabic*]},parsep=0ex]
% \normalsize

\item Meihao Fan, Ju Fan, Yuxin Zhang, \textbf{Shaolei\ Zhang}, Xiaoyong Du, Jie Song, Peng Li, Fuxin Jiang, Tieying Zhang, Jianjun Chen. DeepPrep: An LLM-Powered Agentic System for Autonomous Data Preparation. \textit{Preprint 2026}. [\href{https://arxiv.org/pdf/2602.07371}{PDF}]


\item Haonan Bian, Zhiyuan Yao, Sen Hu, Zishan Xu, \textbf{Shaolei\ Zhang}, Yifu Guo, Ziliang Yang, Xueran Han, Huacan Wang, Ronghao Chen. RealMem: Benchmarking LLMs in Real-World Memory-Driven Interaction. \textit{Preprint 2026} (\textcolor{red2}{CCF-A}). [\href{https://arxiv.org/pdf/2601.06966}{PDF}] [\href{https://github.com/AvatarMemory/RealMemBench}{Code}]

\item Tu Hu, Ronghao Chen, Shuo Zhang, Jianghao Yin, Mou Xiao Feng, Jingping Liu, \textbf{Shaolei\ Zhang}, Wenqi Jiang, Yuqi Fang, Sen Hu, Huacan Wang, Yi Xu. Controlled Self-Evolution for Algorithmic Code Optimization. \textit{Preprint 2026}. [\href{https://arxiv.org/pdf/2601.07348}{PDF}]

\item Xiaomi LLM-Core Team. MiMo-V2-Flash Technical Report. \textit{Technical Report 2026} (\textcolor{red2}{CCF-A}). [\href{https://arxiv.org/pdf/2601.02780}{PDF}]

\item Yizhang Zhu, Liangwei Wang, Chenyu Yang, Xiaotian Lin, Boyan Li, Wei Zhou, Xinyu Liu, Zhangyang Peng, Tianqi Luo, Yu Li, Chengliang Chai, Chong Chen, Shimin Di, Ju Fan, Ji Sun, Nan Tang, Fugee Tsung, Jiannan Wang, Chenglin Wu, Yanwei Xu, \textbf{Shaolei\ Zhang}, Yong Zhang, Xuanhe Zhou, Guoliang Li, Yuyu Luo. A Survey of Data Agents: Emerging Paradigm or Overstated Hype? \textit{Preprint 2025}. [\href{https://arxiv.org/pdf/2510.23587}{PDF}] [\href{https://github.com/HKUSTDial/awesome-data-agents}{List}]


    \item Kangyu Qiao, \textbf{Shaolei\ Zhang}, Yang Feng. IG-Pruning: Input-Guided Block Pruning for Large Language Models. \textit{EMNLP 2025} (\textcolor{red2}{CCF-B}).
    
    \item Mengyu Bu, \textbf{Shaolei\ Zhang}, Zhongjun He, Hua Wu, Yang Feng. AlignX: Advancing Multilingual Large Language Models with Multilingual Representation Alignment. \textit{EMNLP 2025} (\textcolor{red2}{CCF-B}).
    
    \item Shoutao Guo, \textbf{Shaolei\ Zhang}, Zhengrui Ma, Min Zhang, Yang Feng. Agent-SiMT: Agent-assisted Simultaneous Machine Translation with Large Language Models. \textit{TASLP 2025} \textcolor{red1}{(CCF-B)}. [\href{https://arxiv.org/pdf/2406.06910}{PDF}] [\href{https://github.com/ictnlp/SiLLM}{Code}]

    \item Qingkai Fang, Yan Zhou, Shoutao Guo, \textbf{Shaolei\ Zhang}, Yang Feng. LLaMA-Omni2: LLM-based Real-time Spoken Chatbot with Autoregressive Streaming Speech Synthesis. \textit{ACL 2025} (\textcolor{red2}{CCF-A}).

    \item Qingkai Fang, Shoutao Guo, Yan Zhou, Zhengrui Ma, \textbf{Shaolei\ Zhang}, Yang Feng. LLaMA-Omni: Seamless Speech Interaction with Large Language Models. \textit{ICLR 2025}. [\href{https://arxiv.org/pdf/2406.06910}{PDF}] [\href{https://github.com/ictnlp/SiLLM}{Code}]

      
    \item Qingkai Fang, \textbf{Shaolei\ Zhang}, Zhengrui Ma, Min Zhang, Yang Feng. Can We Achieve High-quality Direct Speech-to-Speech Translation Without Parallel Speech Data? \textit{ACL 2024} (\textcolor{red2}{CCF-A}). [\href{https://arxiv.org/pdf/2406.07289}{PDF}] [\href{https://github.com/ictnlp/ComSpeech}{Code}]
    
    \item Shoutao Guo, \textbf{Shaolei\ Zhang}, Yang Feng. Decoder-only Streaming Transformer for Simultaneous Translation. \textit{ACL 2024} (\textcolor{red2}{CCF-A}). [\href{https://arxiv.org/pdf/2406.03878}{PDF}] [\href{https://github.com/ictnlp/DST}{Code}] 

    \item Tian Yu, \textbf{Shaolei\ Zhang}, Yang Feng. Truth-Aware Context Selection: Mitigating the Hallucinations of Large Language Models Being Misled by Untruthful Contexts. \textit{ACL 2024 findings}. [\href{https://arxiv.org/pdf/2403.07556.pdf}{PDF}] [\href{https://github.com/ictnlp/TACS}{Code}]
      
    \item Zhengrui Ma, Qingkai Fang, \textbf{Shaolei\ Zhang}, Shoutao Guo, Min Zhang, Yang Feng. A Non-autoregressive Generation Framework for End-to-End Simultaneous Speech-to-Any Translation. \textit{ACL 2024} (\textcolor{red2}{CCF-A}). [\href{https://arxiv.org/pdf/2406.06937}{PDF}] [\href{https://github.com/ictnlp/NAST-S2x}{Code}] 
      
    \item Shoutao Guo, \textbf{Shaolei\ Zhang}, Yang Feng. Glancing Future for Simultaneous Machine Translation. \textit{ICASSP 2024} (\textcolor{red2}{CCF-B, Oral}). [\href{https://arxiv.org/pdf/2409.06666}{PDF}] [\href{https://github.com/ictnlp/LLaMA-Omni}{Code}]

    \item Zhengrui Ma, \textbf{Shaolei\ Zhang}, Shoutao Guo, Chenze Shao, Min Zhang, Yang Feng. Non-autoregressive Streaming Transformer for Simultaneous Translation. \textit{EMNLP 2022} (\textcolor{red2}{CCF-B, Oral}). [\href{https://aclanthology.org/2023.emnlp-main.314.pdf}{PDF}] [\href{https://github.com/ictnlp/NAST}{Code}]

    \item Shoutao Guo, \textbf{Shaolei\ Zhang}, Yang Feng. Learning Optimal Policy for Simultaneous Machine Translation via Binary Search. \textit{ACL 2023} (\textcolor{red2}{CCF-A}).
    
    \item Shoutao Guo, \textbf{Shaolei\ Zhang}, Yang Feng. Turning Fixed to Adaptive: Integrating Post-Evaluation into Simultaneous Machine Translation. \textit{EMNLP 2022 findings}. [\href{https://arxiv.org/pdf/2210.11900.pdf}{PDF}] [\href{https://github.com/ictnlp/PED-SiMT}{Code}]
    

    \item Tian Yu, \textbf{Shaolei\ Zhang}, Yang Feng. Auto-RAG: Autonomous Retrieval-Augmented Generation for Large Language Models. \textit{Preprint 2024}. [\href{https://arxiv.org/pdf/2411.19443}{PDF}]


    \item Kehao Zhang, \textbf{Shaolei\ Zhang}, Yang Feng. PSO-Merging: Merging Models Based on Particle Swarm Optimization. \textit{Preprint 2025}. [\href{https://arxiv.org/pdf/2508.19839}{PDF}]

    \item Shoutao Guo, \textbf{Shaolei\ Zhang}, Qingkai Fang, Zhengrui Ma, Min Zhang, Yang Feng. FastLongSpeech: Enhancing Large Speech-Language Models for Efficient Long-Speech Processing. \textit{Preprint 2025}. [\href{https://arxiv.org/pdf/2507.14815}{PDF}]

    
\end{enumerate}

% \section{\textcolor{red1}{\faClipboard\ 教学、合作经历}}
% \begin{itemize}[parsep=0.3ex]
% %\fontsize{11pt}{0.2ex}\selectfont
% \large
%     \item 2021中国科学院大学《自然语言处理》教师助教（合作教师：刘洋老师）
%     \item 2022中国科学院大学《自然语言处理》教师助教（合作教师：冯洋老师）
%     \item 百聆大模型作为基座模型服务于国科大实践课《通用大模型原理及训练实践》 
%     \item 与组内博士硕士（博四、博三、博二、研二）合作发表国际会议论文8篇
% \end{itemize}


\end{document}
